\documentclass[12pt]{article}
\usepackage[utf8]{inputenc}
\usepackage[T1]{fontenc}
\usepackage{hyperref}
\title{Tebul Sign Language Language Resources for K-12 Curriculum Development}
\author{Compiled by Claude Sonnet 5 (claude-sonnet-5) for Andr\'as Kornai}
\date{}
\begin{document}
\maketitle

\section{Introduction}
Tebul Sign Language (ISO 639-3 tsy) is a village sign language of Uluban, Mali, in the Dogon region. Its users are mostly speakers of Tebul Dogon. Sources report about 500 users in 2007 and a high rate of deafness in the community. Nearly all residents may know some signs. Very little is documented.

\section{Primary Educational Resources}
No curricula, primers or schools for Tebul Sign Language were found in search results.

\section{Linguistic and Reference Materials}
\begin{itemize}
\item Wikipedia article on Tebul Sign Language: \url{https://en.wikipedia.org/wiki/Tebul_Sign_Language}
\item Wikipedia article on Tebul Dogon, the spoken language of the community: \url{https://en.wikipedia.org/wiki/Tebul_Dogon}
\item Wikipedia article on Douentza Sign Language, a different sign language of the region: \url{https://en.wikipedia.org/wiki/Douentza_Sign_Language}
\end{itemize}

\section{Religious and Community Resources}
None found.

\section{Digital and Audio Resources}
\begin{itemize}
\item Information from Roger Blench (2012), ``Tebul Ure, a language of the Dogon group in Northern Mali and its affinities'', as shown on endangeredlanguages.com. It mentions narrative videos recorded with translations into Tebul Ure and French. I did not confirm that the videos are available online: \url{https://endangeredlanguages.com/lang/8735?hl=de}
\end{itemize}

\section{Curriculum Development Context}
The evidence is thin. A curriculum would first need documentation led by Deaf signers in Uluban, and it should not borrow from Douentza Sign Language without evidence that the two are related.

\end{document}